% 结构版本 v0.2：用户已确认审阅稿（D015）。5.1–5.3 定稿（V013，2026-09-26 用户对话审定＋两轮 GPT 评审逐条仲裁）；5.4–5.5 仍为骨架（标题未随 D038 大纲随迁，归整合窗）。
% V013 起草核料：官方评估器 P2 语义行号级核验（Task 构造/同核零边界搬运/按核判定/500 cycles 释放下界/计费 added=partition+spill/step2 容量与换出失败条件/_prioritize_task_seq 重排）＋题面 1.5、问题 2、1.4、B.4、C.1/C.2、D.2/D.5 原文＋基准写法提取（2025A 第 5 章＝机制全新全量重定义型，本文为场景适配型，范式对齐 2023A 5.2"骨架沿用＋差异声明＋结构改造"）。
% V013 与 W-C 窗口（V012/A049，4.1–4.2）符号协调：本文件不自造符号，沿用第四章已定记号——决策 D=(x,{π_k})、核心分配 φ（序列归属导出）、标签集合 S 与子图 O_s、子图依赖图 Q=(S,E_dep)、方案级合法集合 𝒟plan（4.2.3 五条静态约束）、场景 B 侧对应 mk_B/add_B/𝒟feas^B；总模型 Φ_B 的 q_B 一般定义（完成有效评估指示）与 𝒟feas^B 成员等价，为第六章同型映射留接口。预览稿自造的 a/τ_j/Ω/x(方案整体) 全部废弃。
% V013 仲裁要点（两轮 GPT 评审＋用户终裁，逐条见 09/A050）：
% ①容量失败条件取题面 C.2 完整两类（"单个张量超过容量或不存在合法换出对象"）——代码路径合流（step2 无前置单张量检查）不等于题面语义，不以实现路径删语义；
% ②B_cross 集合化 P(t)×C(t)（与评估器 producer_cores×consumer_cores 同构）；题面无单生产者保证，不设单生产者退化句；
% ③场景 B＝两项并列机制（核级 Task 组织＋跨核同步信号），撤"全部差异源于一条规则"单一因果（段1/段5 同框架）；
% ④驻留窗口等执行层量用"影响"不用"决定"（方案可控结构层 vs 评估派生执行层两分自洽）；
% ⑤"不因子图边界新增搬运""固定 Task 等待被取消"限定表述（同核 spill 仍产生搬运；依赖/管道/带宽等待不属"固定等待"）；
% ⑥评价映射总模型 Φ_B（x∈Ω 提交层合法、q_B=1 执行层可行的两层合法性，与 1.3.1 V010 分层同构，为第六章同型映射留接口）；
% ⑦B_B 不进入接受准则（与 v5 实现 mk2<best_mk 及 1.3.3 V010"接受准则只看总体任务执行时间"一致；不虚构加权/词典序）；
% ⑧同步释放时刻 r_sync 与实际发射时刻 s_in 分离（500 cycles 是释放下界，实际发射由完整模拟确定）；
% ⑨5.3.4 五条长句改四列对照表（表 5-1＝TBL-013，registry/05 已同轮登记；与 TBL-002 分工待 T024 定夺——本表管 P1→P2 模型维度对照，TBL-002 若留表应收窄三问决策边界全览）。
% 延后项：场景 A 等待常量（100/1000）数值与完整机制已落 4.1（V012/W-C），本文件 \ref{sec:q1-desc} 回指；C_inner/C_cross 不引入符号（5.2/5.3 不再用到，按 1.2.3 注记"确有使用时引入"以文字转述）；F09/算法 2 落位 5.3 节尾，依赖 4.3 定稿＋#7 方法冻结，引用句与 figure/algorithm 环境随图定版补入（同 2.5 F01 先例）。
% 符号首现（待 3.2/TBL-003 收录；D/x/π_k/φ/S/O_s/Q/E_dep/𝒟plan 沿用 4.2.1/4.2.3，不重复登记）：G^p(t)、H(t)、R、X、P(t)、C(t)、L_t、B_cross、B_reread、B_spill、δ_sync、e^out/r^sync/s^in、T_input、Φ_B、mk_B、add_B、q_B、𝒟feas^B。
\section{场景 B 下的缓存感知多核调度模型}\label{sec:q2}

\subsection{场景 B 与场景 A 的差异}
% 状态：V013 定稿（2026-09-26 用户对话审定＋两轮 GPT 评审仲裁）。场景 B 机制事实全章唯一展开处（1.5.2 定义层一句、2.3 定性权衡、1.2.3 Task 组织均回指此处）。素材：题面 1.5/问题 2/附录 B.4/C.1/C.2/D.2＋评估器核验；无实验数据、无方法名、不挂文献。
问题二与问题一共享同一提交层决策空间与正式优化目标（见 1.5.2 节），二者的差异不在方案表示，而在场景 B 的执行与评价语义。首先是 Task 的组织方式：场景 A 无核间同步机制，每个子图独立构成一个 Task，Task 之间的数据传递一律经 DDR 中转；场景 B 在官方评估规则下，将同一核心上被调度的全部子图合并为一个 Task。子图在核内的执行顺序仍由子图调度方案给出；对合并后的 Task，官方核内调度在不破坏数据依赖的前提下，按方案给定的同核子图顺序调整核内操作的执行次序。因此，切图与调度方案的表示和基本合法性条件（见 1.2.3 节）在两个场景中完全一致，场景差异不进入方案本身，而进入方案之下的执行与评价过程。

合并 Task 的直接后果，是同核子图边界的消隐。同核的生产—消费依赖（无论两端是否属于同一子图）不再插入边界搬运：前序操作的输出张量驻留于本核 L1/UB，供后续子图直接使用；对来自 DDR 的图输入，同一核心只搬入一次，核内的全部子图共享同一副本。由此，场景 B 中复用是否可能，由方案的静态结构决定——生产方与消费方是否位于同一核心；而复用是否实际发生，还取决于容量条件——张量须持续驻留至其再次使用，且期间未被缓存换入换出（spill）移出。前者是算法可以直接控制的，后者由官方评估程序的核内调度在容量约束下确定（见 1.3.3 节）。

跨核依赖仍经 DDR 中转，但组织方式与场景 A 不同：对每条生产方与消费方分属不同核心的依赖，源核心插入 COPY\_OUT 将数据写入 DDR，目标核心插入 COPY\_IN 读回，两者传递同一 DDR 张量；目标核心的 COPY\_IN 须待源核心 COPY\_OUT 完成并经过固定的同步延迟（500 cycles）后方可发射。与场景 A 的差异体现在判定粒度与等待结构上：场景 A 按子图（Task）边界判定搬运，相邻 Task 即使位于同一核心也须经 DDR 中转，且 Task 激活受同核切换等待与跨核前驱等待约束（其机制与数值见~\ref{sec:q1-desc} 节）；场景 B 按核心判定搬运，同核的跨子图依赖不因子图边界新增搬运，场景特有的固定 Task 等待被取消，跨核通信的固定同步开销体现为每条跨核通路的 500 cycles 延迟。题面以“同核单位通信成本低于跨核”概括这一不对称；在官方评估语义下，其具体形式即同核通路免除边界搬运与固定等待、跨核通路计入往返搬运与同步延迟。

同核驻留的代价由容量承担。跨越多个子图的驻留数据持续占用核内 L1/UB（两类缓存分别检查容量，见 1.2.1 节），可能压缩后续子图执行时的可用空间；当核内调度的驻留需求超出容量时，官方评估程序自动插入缓存换出与换入，其搬运计入总额外数据搬运量并消耗 DDR 带宽；若单个张量超过容量或不存在合法换出对象，则该 Task 评估失败。因此，题面所述“同一核内的子图调度顺序影响依赖等待、张量驻留时间和缓存换入换出”在场景 B 中完整生效：同核子图的先后顺序影响驻留张量在核内的存续时长与并发驻留规模，进而影响容量压力与换入换出的规模。相对场景 A 中同核顺序主要影响 Task 激活时序（见 2.2 节），场景 B 赋予同核顺序第二重作用。

综上，场景 B 的评价语义变化可归结为两项机制的叠加：其一，Task 由子图级聚合到核级，使同核的生产—消费依赖能够驻留复用，也使核内容量压力由单个子图的工作集扩大为同核子图与驻留数据的并发集合；其二，跨核依赖引入同步信号传递，使跨核通路在 DDR 往返搬运之外承担固定的同步延迟。前者划定哪些数据能够留在核内，后者规定离开核心的数据如何返回；核心分配在划分负载的同时，为每条跨子图依赖选定其中一条通路。5.2 节将上述机制整理为驻留与通信的形式化表达，5.3 节据此给出场景 B 的调度优化模型。

\subsection{核内数据驻留与通信成本模型}
% 状态：V013 定稿。内容边界（D048）：机制规则引用 5.1 不复述，本节只建状态、关系与成本表达；记号沿用 4.2.1/4.2.3（D=(x,{π_k})、φ、S、Q=(S,E_dep)、𝒟plan），本节不重复定义。
给定一份决策 $D\in\mathcal{D}_{\mathrm{plan}}$（记号见 4.2.1、4.2.3 节），场景 B 的执行语义由两级派生结构完全确定：其一为 Task 划分，其二为由 Task 划分导出的逐张量数据通路。本节对二者建立形式化表达，并给出场景 B 下通信成本的构成；该表达不改变方案的提交形式，用于后文模型的约束与代价分析。

\subsubsection{同核数据复用与驻留条件}
沿用 4.2.1 节的记号：切图映射 $x$ 给出标签集合 $S$ 及各子图 $O_s$，调度序列 $\{\pi_k\}$ 给出核心分配 $\phi$ 与各核子图执行顺序，子图间依赖构成子图依赖图 $Q=(S,\,E_{\mathrm{dep}})$（1.2.3、4.2.3 节）。场景 B 的 Task 划分即由 $\{\pi_k\}$ 给出：第 $k$ 个 Task 恰为核 $k$ 上按 $\pi_k$ 顺序执行的全部子图，Task 划分与核心一一对应。

对每个由核内操作生产的张量 $t$，记其生产子图集合为 $G^{\mathrm{p}}(t)\subseteq S$、消费子图集合为 $\mathcal{H}(t)\subseteq S$。每一条跨子图生产—消费关系 $(s,s',t)$（$s\in G^{\mathrm{p}}(t)$、$s'\in\mathcal{H}(t)$、$s\neq s'$）被核心分配划分为两类：
\[ \mathcal{R}=\{(s,s',t):\ \phi(s)=\phi(s')\},\qquad \mathcal{X}=\{(s,s',t):\ \phi(s)\neq\phi(s')\}. \]
$\mathcal{R}$ 中的关系构成\textbf{潜在驻留复用对}：张量 $t$ 自生产完成起驻留于核 $\phi(s)$ 的 L1/UB，供 $s'$ 直接使用，不产生边界搬运（官方机制给定的规则，见 5.1 节）；$\mathcal{X}$ 中的关系构成\textbf{跨核通信对}，其通路与计费见 5.2.2 节。图输入张量位于 DDR、无核内生产方，其搬入方式与代价归入 5.2.3 节的重复读取项。

潜在复用是否转化为实际复用，取决于容量条件：$t$ 须驻留至其消费操作执行，且期间不因容量不足被换出。驻留窗口自 $t$ 生产完成起、至其在核内最后一次使用止，窗口长度受 $\pi_k$ 中生产、消费子图相对次序的影响，其实际边界还取决于 Task 内操作依赖与核内调度的结果；窗口内的并发驻留集合与 L1/UB 容量的关系决定换入换出是否发生（规则见 5.1 节与附录 C）。据此，本文区分两个层次：结构层面的 $\mathcal{R}$ 可由方案静态计算并用于算法构造；执行层面的实际驻留由官方评估确定，并通过评价结果体现，与 1.3.3 节的容量口径一致。

\subsubsection{跨核数据搬运与同步}
场景 B 下，跨核数据通路按（源核、目标核、张量）组织。对由核内操作生产的张量 $t$，记其生产核集合与消费核集合为
\[ P(t)=\{\phi(s):\ s\in G^{\mathrm{p}}(t)\},\qquad C(t)=\{\phi(s'):\ s'\in\mathcal{H}(t)\}, \]
每个满足 $i\neq j$ 的核对 $(i,j)\in P(t)\times C(t)$ 对应一条跨核通路：源核 $i$ 的 COPY\_OUT 写入 DDR、目标核 $j$ 的 COPY\_IN 读回，传递同一 DDR 张量，搬运字节数各为 $s_t$；同一目标核上的多个消费子图共享该次 COPY\_IN。

同步约束表达为：对每条通路 $\ell$，其同步释放时刻为源端 COPY\_OUT 完成时刻加固定同步延迟 $\delta_{\mathrm{sync}}=500$ cycles，
\[ r^{\mathrm{sync}}(\ell)=e^{\mathrm{out}}(\ell)+\delta_{\mathrm{sync}}, \]
目标 COPY\_IN 的实际发射时刻不早于同步释放时刻：
\[ s^{\mathrm{in}}(\ell)\ \ge\ r^{\mathrm{sync}}(\ell). \]
其中 $e^{\mathrm{out}}(\ell)$ 由源核 COPY\_OUT 在共享带宽下的完成事件确定，$s^{\mathrm{in}}(\ell)$ 由完整多核模拟在依赖、管道与带宽约束下共同确定（多核模拟的口径见 1.4 节）；$\delta_{\mathrm{sync}}$ 与数据量无关，只取决于通路的跨核属性。记 $\mathcal{L}_t=\{(i,j):\ i\in P(t),\ j\in C(t),\ i\neq j\}$，则中间张量的跨核边界搬运量为
\[ B_{\mathrm{cross}}=2\sum_{t} s_t\,\bigl|\mathcal{L}_t\bigr|, \]
即每条通路计入写入与读取各一次。$B_{\mathrm{cross}}$ 完全由方案静态决定，是 1.3.2 节跨组生产—消费张量字节特征在场景 B 下的细化。

\subsubsection{通信成本、容量压力与 spill}
场景 B 下总额外数据搬运量的构成仍为 1.4 节所述三类，计费口径随场景规则变化。其一为跨核边界搬运，即式 $B_{\mathrm{cross}}$。其二为重复读取：题设计算图的每个 DDR 输入只搬入一次，而场景 B 中图输入由各消费核独立搬入，消费核数为 $k(t)$ 的图输入张量计入 $k(t)$ 次搬入，超出原图一次的部分为
\[ B_{\mathrm{reread}}=\sum_{t\in T_{\mathrm{input}}} \bigl(k(t)-1\bigr)\,s_t, \]
其中 $T_{\mathrm{input}}$ 为图输入张量集合。其三为缓存换入换出 $B_{\mathrm{spill}}$：驻留需求超出容量时核内调度自动插入换出与换入，其搬运计入该项；$B_{\mathrm{spill}}$ 依赖核内调度的具体决策，无法由方案静态算得，本文不以近似值替代（见 1.3.3 节）。题设保证每个最终结果只搬出一次，图输出的写回沿用原图口径，不在此口径下产生新增项；三类之和即总额外数据搬运量，由官方评估程序统计。

容量压力方面，驻留张量按其位置属性分别计入 L1 与 UB 的占用（两类容量独立，数值见表~\ref{tab:hw-config}），驻留窗口内未释放的张量构成并发驻留集合，其规模与构成受 Task 内子图构成及其顺序的共同影响（5.2.1 节）。需要强调，$\mathcal{R}$ 中全部张量大小之和只刻画潜在驻留的总规模，既非任一时刻的实际占用，也非必须同时驻留的集合；实际驻留峰值由官方核内调度在评估中确定。因此，本文不将其简单相加作为容量约束，容量对方案的作用经由 $B_{\mathrm{spill}}$ 与执行时间进入评价。上述表达中，$\mathcal{R}$、$\mathcal{X}$ 与 $B_{\mathrm{cross}}$、$B_{\mathrm{reread}}$ 为静态可算量，构成第五章后续算法构造与代价估计的直接依据；$B_{\mathrm{spill}}$ 与实际驻留经由评价获得。

\subsection{缓存约束下的多核调度模型}
% 状态：V013 定稿。评价映射总模型：与 4.2.2 的两级可行域同型对齐——𝒟plan 沿用 4.2.3（五条静态约束场景不变），场景 B 可行集合 𝒟feas^B；q_B 一般定义（完成有效评估指示）与 𝒟feas^B 成员等价，两层合法性与 1.3.1（V010）分层同构，为第六章 Φ_C 同型映射留接口。
在 5.1 节的机制差异与 5.2 节的形式化表达基础上，本节给出场景 B 下的多核切图与调度优化模型。本题的提交方案并不直接决定核内换入换出与多核执行时间线，这些量由官方评估规则从提交方案确定性派生（见 1.2.3、2.5 节）。因此，本文将问题二建模为\textbf{合法方案空间上的场景 B 执行映射及其搜索问题}：沿用 4.2.1 节的决策 $D=(x,\{\pi_k\})$ 与 4.2.3 节的方案级合法集合 $\mathcal{D}_{\mathrm{plan}}$——五条静态条件只涉及方案表示与依赖结构，不随场景改变。场景 B 的官方执行过程确定评价映射
\[ \Phi_B:\ D\ \mapsto\ \bigl(\mathrm{mk}_B(D),\ \mathrm{add}_B(D),\ q_B(D)\bigr), \]
其中 $\mathrm{mk}_B(D)$ 为场景 B 语义下多核模拟输出的总体任务执行时间，$\mathrm{add}_B(D)$ 为总额外数据搬运量（构成见 5.2.3 节），$q_B(D)$ 指示方案能否在场景 B 执行语义下完成有效评估；$q_B(D)=1$ 的方案构成场景 B 可行集合 $\mathcal{D}_{\mathrm{feas}}^{B}\subseteq\mathcal{D}_{\mathrm{plan}}$。问题二即
\begin{equation}\label{eq:q2-object}
\min_{D\in\mathcal{D}_{\mathrm{feas}}^{B}}\ \mathrm{mk}_B(D),
\end{equation}
并以 $\mathrm{add}_B(D)$ 作为辅助评价指标，用于衡量方案的通信代价（其角色的进一步说明见 5.3.2 节）。以下 5.3.1 节说明决策空间的组成与含义变化，5.3.2 节说明评价方式，5.3.3 节说明 $q_B$ 的场景 B 特有内容，5.3.4 节以变化清单归纳相对问题一的模型差异。

\subsubsection{决策变量与共用定义}
决策变量沿用问题一：$D=(x,\{\pi_k\})$ 与两个方案字段一一对应（4.2.1 节），核心分配 $\phi$ 由序列归属导出。场景 B 不新增任何提交层面的决策变量；5.2 节的 Task 划分、潜在复用对集合 $\mathcal{R}$、跨核通信对集合 $\mathcal{X}$ 与通路集合 $\{\mathcal{L}_t\}$ 均为方案导出的派生量。

决策变量的含义在场景 B 下发生实质变化。其一，核心分配在负载划分之外同时决定每条跨子图依赖的数据通路：$\phi$ 的取值直接划定 $\mathcal{R}$ 与 $\mathcal{X}$ 的边界，即为该依赖选定同核驻留复用或跨核中转（5.2.1 节）。其二，同核子图顺序除保持依赖一致性外，经核内调度的顺序重排（5.1 节）影响驻留窗口长度与并发驻留规模，进而影响容量压力与换入换出（5.2.1、5.2.3 节）。两者叠加，使切图粒度与分核方式共同决定潜在可复用空间，复用收益又受容量条件制约，问题二的决策耦合较问题一更紧。

\subsubsection{目标函数与评价方式}
优化目标沿用问题一：以总体任务执行时间 $\mathrm{mk}_B(D)$ 最小为正式优化目标，总额外数据搬运量 $\mathrm{add}_B(D)$ 为辅助评价指标（二者定义见 1.4 节）。与实际求解一致，本文的最终接受准则仅以总体任务执行时间 $\mathrm{mk}_B(D)$ 的严格改进为依据；$\mathrm{add}_B(D)$ 不进入接受准则，用于结果分析与方案通信代价的衡量——这对应题面“以缩短总体任务执行时间为主要目标、兼顾总额外数据搬运量”的要求。目标函数形式与问题一相同（式~\ref{eq:q1-object}），变化在于评价环境：场景 B 的官方评估按 5.1 节机制构造 Task 与边界搬运、执行核内调度（含顺序重排、容量检查与换入换出），并在共享 DDR 带宽与跨核同步延迟下完成多核模拟，$\mathrm{mk}_B$ 与 $\mathrm{add}_B$ 均在该环境下由官方评估程序给出。求解过程中如需利用 5.2 节静态可算量的内部代价与加速评估实现，其构造方式见 5.4 节。

\subsubsection{缓存容量、依赖及其他约束}
$\mathcal{D}_{\mathrm{plan}}$ 所含的五条静态条件（4.2.3 节）保证方案在提交层面合法；$q_B$ 在此之上要求方案在场景 B 执行语义下完成有效评估，执行层检查由官方评估流程承担。其中场景 B 特有的关键项是 L1/UB 容量的处理：题面要求严格满足容量约束，在官方评估语义下，该约束不由方案静态校验，而由核内调度在评估中强制执行——驻留需求超出容量时自动插入换入换出予以化解，单个张量超过容量或不存在合法换出对象时 Task 评估失败（5.1 节）。据此，本文对容量的处理是：不设置静态容量约束或容量近似（1.3.3 节），将容量压力作为方案的评价后果，经由换入换出搬运与执行时间体现。

\subsubsection{相对问题一的模型变化与合法性检查}
% 表 TBL-013（表 5-1）：V013 首次成表，05/registry 已同轮登记（D029）；与 TBL-002（2.5，去留未定）分工待 T024 定夺——本表管 P1→P2 模型维度对照，TBL-002 若留表应收窄为三问决策边界全览。
以表~\ref{tab:q2-diff}归纳场景 B 相对问题一的模型差异。

\begin{table}[htbp]
  \centering
  \caption{场景 B 相对问题一的模型变化对照}
  \label{tab:q2-diff}
  \begin{tabular}{lp{2.3cm}p{2.9cm}p{3.7cm}}
    \toprule
    模型维度 & 场景 A & 场景 B & 对优化的影响 \\
    \midrule
    Task 粒度 & 每个子图一个 Task & 每核一个 Task & 同核跨子图复用成为可能 \\
    边界搬运判定 & 按子图（Task）边界 & 按核心 & 核心分配决定每条依赖的数据通路 \\
    固定等待 & Task 激活的同核切换与跨核前驱等待 & 逐跨核通路同步延迟（500 cycles） & 等待由 Task 级转为通路级 \\
    核内工作集 & 单个子图 & 同核子图与驻留数据并发 & 同核顺序影响驻留时长与换入换出 \\
    新增搬运构成 & 跨 Task 边界＋按 Task 重复读取＋换入换出 & 跨核边界＋按核重复读取＋换入换出 & 计费口径见 5.2.2、5.2.3 节 \\
    \bottomrule
  \end{tabular}
  \begin{minipage}{0.9\linewidth}
    \small 注：本表据 5.1--5.2.3 节归纳；场景 A 侧机制与数值见~\ref{sec:q1-desc} 节。
  \end{minipage}
\end{table}

因此，问题二未改变提交层决策变量与目标函数，改变的是同一方案对应的执行映射与成本结构；第四章的求解框架可整体沿用，针对上述差异的适配方式见 5.4 节。
% TODO(F09/算法2, 4.3 定稿＋#7 方法冻结后补)：F09 子图重构与 M/V 双管道交错机制示意（双子图，标注非实测时间线）引用句＋figure 环境；算法 2 重标伪代码 algorithm 环境。落点 5.3 节尾（大纲现行落位；最终落点 5.3 vs 随 4.3 重标内容归第四章，随 T024 定夺）。

\subsection{场景 B 下的求解算法}
\pending{内容 TODO：依据题面、最终采用的实现及适用证据补充本节；具体算法、公式和结果尚未填写。}

\subsubsection{算法总体思想与选择依据}
\pending{内容 TODO：依据题面、最终采用的实现及适用证据补充本节；具体算法、公式和结果尚未填写。}

\subsubsection{相对问题一的适配策略}
\pending{内容 TODO：依据题面、最终采用的实现及适用证据补充本节；具体算法、公式和结果尚未填写。}

\subsubsection{算法流程、伪代码与停止条件}
\pending{内容 TODO：依据题面、最终采用的实现及适用证据补充本节；具体算法、公式和结果尚未填写。}

\subsubsection{算法复杂度与计算预算}
\pending{算法名称及改动内容 TODO；不预设场景 B 方法一定是问题一算法的局部改进。}

\subsection{实验结果与分析}
\pending{内容 TODO：依据题面、最终采用的实现及适用证据补充本节；具体算法、公式和结果尚未填写。}

\subsubsection{实验设置与场景对照协议}
\pending{引用公共设置，补充场景 B 特有配置和对照设计。}

\subsubsection{1～5 核平均加速比}
\pending{内容 TODO：依据题面、最终采用的实现及适用证据补充本节；具体算法、公式和结果尚未填写。}

\subsubsection{Makespan 与总额外数据搬运量}
\pending{内容 TODO：依据题面、最终采用的实现及适用证据补充本节；具体算法、公式和结果尚未填写。}

\subsubsection{场景 A 与场景 B 对比}
\pending{区分固定合法方案的场景对照与各场景分别优化后的算法对照。}

\subsubsection{缓存复用效果与关键策略消融}
\pending{内容 TODO：依据题面、最终采用的实现及适用证据补充本节；具体算法、公式和结果尚未填写。}

\subsubsection{求解成本与典型案例}
\pending{内容 TODO：依据题面、最终采用的实现及适用证据补充本节；具体算法、公式和结果尚未填写。}

\subsubsection{算法结果小结}
\pending{内容 TODO：依据题面、最终采用的实现及适用证据补充本节；具体算法、公式和结果尚未填写。}
